% !TEX root = ../main.tex
\begin{ouverture}[forceacide]{9}{Force des acides et des bases}{Partie 1 · Constitution et transformations de la matière}
{Produit ionique de l'eau\enspace·\enspace Constante d'acidité\enspace·\enspace Diagrammes de prédominance\enspace·\enspace Solutions tampons}
\objectif{Utiliser le produit ionique de l'eau pour relier pH et concentration en ions hydroxyde.}
\objectif{Calculer le pH d'une solution d'acide fort ou de base forte.}
\objectif{Associer la constante d'acidité d'un couple à la force de l'acide et au taux d'avancement final.}
\objectif{Exploiter un diagramme de prédominance ou de distribution ; comprendre un indicateur coloré.}
\objectif{Décrire le comportement d'une solution tampon.}
\prerequis{Transformations acide-base, couples acide-base (chapitre~1).}
\prerequis{Quotient de réaction, constante d'équilibre, taux d'avancement final (chapitre~7).}
\prerequis{Fonction logarithme décimal ; équation du second degré.}
\end{ouverture}
\essentiel
\essSection{1}{Le produit ionique de l'eau}
\begin{definition}[Autoprotolyse et produit ionique]
L'eau réagit avec elle-même selon l'équilibre d'\textbf{autoprotolyse}
\[\ce{2 H2O(l) <=> H3O^+(aq) + HO^-(aq)}.\]
La constante d'équilibre de cette réaction est le \textbf{produit ionique de l'eau} :
\[K_\mathrm{e}=\frac{[\ce{H3O+}]_\text{éq}\,[\ce{HO-}]_\text{éq}}{(\cst)^2},\qquad K_\mathrm{e}=\num{1.0e-14}\ \text{à 25~°C},\qquad
\mathrm{p}K_\mathrm{e}=-\log K_\mathrm{e}=14{,}0.\]
$K_\mathrm{e}$ ne dépend que de la température. Dans toute solution aqueuse, $\pH+\mathrm{pOH}=\mathrm{p}K_\mathrm{e}$, avec
$\mathrm{pOH}=-\log\left([\ce{HO-}]/\cst\right)$.
\end{definition}
\begin{propriete}[Solutions acides, neutres, basiques]
Une solution est \textbf{neutre} si $[\ce{H3O+}]=[\ce{HO-}]$, c'est-à-dire si $\pH=\frac12\,\mathrm{p}K_\mathrm{e}$, soit 7,0
à 25~°C. Elle est \textbf{acide} si $\pH<7{,}0$ et \textbf{basique} si $\pH>7{,}0$, toujours à 25~°C : à une autre
température, la frontière se déplace avec $K_\mathrm{e}$.
\end{propriete}
\begin{savoirfaire}[Passer du pH à la concentration en ions hydroxyde]
Calculer $[\ce{H3O+}]=\cst\times10^{-\pH}$, puis utiliser le produit ionique :
$[\ce{HO-}]=\dfrac{K_\mathrm{e}\,(\cst)^2}{[\ce{H3O+}]}$. Par exemple, à $\pH=11{,}0$ et à 25~°C :
$[\ce{H3O+}]=\qty{1,0e-11}{mol.L^{-1}}$ et $[\ce{HO-}]=\qty{1,0e-3}{mol.L^{-1}}$.
\end{savoirfaire}
\essSection{2}{Acides forts et bases fortes}
\begin{definition}[Acide fort, base forte]
Un \textbf{acide fort} réagit totalement avec l'eau : $\ce{AH + H2O -> A- + H3O+}$. Une \textbf{base forte} réagit
totalement : $\ce{B + H2O -> BH+ + HO-}$. Le chlorure d'hydrogène \ce{HCl} (acide chlorhydrique) est un acide fort ;
l'ion hydroxyde apporté par la soude \ce{NaOH} est une base forte.
\end{definition}
\begin{propriete}[pH des solutions d'acide fort et de base forte]
Pour une concentration apportée $c$ ni trop grande ni trop faible (entre \qty{1e-6}{mol.L^{-1}} et \qty{5e-2}{mol.L^{-1}}
environ), à 25~°C :
\[\text{acide fort : }\pH=-\log\frac{c}{\cst}\qquad\text{et}\qquad\text{base forte : }\pH=14{,}0+\log\frac{c}{\cst}.\]
\end{propriete}
\begin{exemple}
Une solution d'acide chlorhydrique à $c=\qty{2,0e-3}{mol.L^{-1}}$ a pour pH $-\log(\num{2.0e-3})=2{,}70$. Une solution
de soude à $c=\qty{1,0e-3}{mol.L^{-1}}$ a pour pH $14{,}0+\log(\num{1.0e-3})=11{,}0$.
\end{exemple}
\begin{attention}
Ces relations supposent que l'acide ou la base apporté impose le pH. Pour une solution très diluée
($c\lesssim\qty{1e-6}{mol.L^{-1}}$), l'autoprotolyse de l'eau n'est plus négligeable : le pH tend vers 7,0, il ne suit
plus $-\log(c/\cst)$ (défi D09).
\end{attention}
\essSection{3}{La constante d'acidité}
\begin{definition}[Constante d'acidité]
La réaction d'un acide AH avec l'eau, $\ce{AH(aq) + H2O(l) <=> A^-(aq) + H3O^+(aq)}$, a pour constante d'équilibre la
\textbf{constante d'acidité} du couple $\mathrm{AH}/\mathrm{A}^-$ :
\[K_\mathrm{A}=\frac{[\ce{A-}]_\text{éq}\,[\ce{H3O+}]_\text{éq}}{[\ce{AH}]_\text{éq}\,\cst},\qquad \pKa=-\log K_\mathrm{A}.\]
Un acide est d'autant plus \textbf{fort} (dissocié) que $K_\mathrm{A}$ est grand, c'est-à-dire que son $\pKa$ est petit ;
à concentration égale, son taux d'avancement final est alors plus élevé et le pH plus faible.
\end{definition}
\begin{propriete}[Quelques pKa à 25 °C]
\par\smallskip
{\small\renewcommand{\arraystretch}{1.25}
\begin{tabularx}{\linewidth}{@{}lYc@{}}\toprule
\otTete{Couple} & \otTete{Acide présent dans} & \otTete{$\pKa$}\\\midrule
\ce{HCOOH}/\ce{HCOO-} & le venin de fourmi & 3,75\\
\ce{C6H5COOH}/\ce{C6H5COO-} & les conservateurs E210 à E213 & 4,20\\
\ce{CH3COOH}/\ce{CH3COO-} & le vinaigre & 4,76\\
\ce{CO2,H2O}/\ce{HCO3-} & les boissons gazeuses & 6,3\\
\ce{HClO}/\ce{ClO-} & l'eau des piscines & 7,5\\
\ce{NH4+}/\ce{NH3} & les engrais & 9,2\\\bottomrule
\end{tabularx}}
\end{propriete}
\begin{propriete}[Relation entre pH et pKa]
En prenant le logarithme de l'expression de $K_\mathrm{A}$, on obtient, dans toute solution contenant le couple :
\[\pH=\pKa+\log\frac{[\ce{A-}]_\text{éq}}{[\ce{AH}]_\text{éq}}.\]
En particulier, $\pH=\pKa$ lorsque $[\ce{A-}]_\text{éq}=[\ce{AH}]_\text{éq}$.
\end{propriete}
\begin{savoirfaire}[Calculer le pH d'une solution d'acide faible]
Dresser le tableau d'avancement de la réaction avec l'eau pour une concentration apportée $c$. À l'équilibre,
$[\ce{A-}]=[\ce{H3O+}]=h$ et $[\ce{AH}]=c-h$ : écrire $K_\mathrm{A}=\dfrac{h^2}{c-h}$, puis résoudre l'équation du second
degré $h^2+K_\mathrm{A}h-K_\mathrm{A}c=0$ et ne garder que la racine positive.
\end{savoirfaire}
\begin{exemple}
Pour l'acide éthanoïque ($K_\mathrm{A}=\num{1.7e-5}$) à $c=\qty{1,0e-2}{mol.L^{-1}}$, la racine positive est
$h=\qty{4,0e-4}{mol.L^{-1}}$, d'où $\pH=3{,}39$ : on retrouve l'avancement calculé par dichotomie au chapitre~7.
\end{exemple}
\essSection{4}{Prédominance, indicateurs et tampons}
\begin{propriete}[Diagramme de prédominance]
L'espèce qui \textbf{prédomine} est celle dont la concentration est la plus grande. D'après la relation
$\pH=\pKa+\log\left([\ce{A-}]/[\ce{AH}]\right)$ : si $\pH<\pKa$, la forme acide AH prédomine ; si $\pH>\pKa$, la forme
basique \ce{A-} prédomine.
\end{propriete}
\schema{09-predominance}
\begin{codecote}[0.58]
\begin{pythoncol}
import numpy as np
import matplotlib.pyplot as plt
pKa = 4.76     # acide éthanoïque
pH = np.linspace(2, 8, 300)
aB = 1 / (1 + 10**(pKa - pH))
plt.plot(pH, 1 - aB, label="CH3COOH")
plt.plot(pH, aB, label="CH3COO-")
plt.xlabel("pH")
plt.legend()
plt.show()
\end{pythoncol}
\cote
\fontsize{9.5}{12.5}\selectfont
\capacite{Capacité numérique}\par\smallskip
La proportion de la forme basique vaut
$\alpha=\dfrac{[\ce{A-}]}{[\ce{A-}]+[\ce{AH}]}=\dfrac{1}{1+10^{\,\pKa-\pH}}$, d'après la relation entre pH et $\pKa$.
Le programme trace les deux \textbf{diagrammes de distribution} du couple : les courbes se croisent à $\pH=\pKa$.
\end{codecote}
\schema{09-python-distribution}
\begin{propriete}[Indicateurs colorés]
Un indicateur coloré acide-base est un couple $\mathrm{HIn}/\mathrm{In}^-$ dont les deux formes ont des couleurs
différentes. La couleur observée est celle de la forme qui prédomine nettement ; autour de $\pKa$, les deux formes
coexistent et les couleurs se mélangent : c'est la \textbf{zone de virage}, large d'environ deux unités de pH.
\end{propriete}
\begin{definition}[Solution tampon]
Une \textbf{solution tampon} contient un acide faible et sa base conjuguée en quantités comparables. Son pH, voisin du
$\pKa$ du couple, \textbf{varie peu} par ajout modéré d'acide ou de base, ou par dilution modérée. Le sang est tamponné
à $\pH=7{,}4$ par le couple \ce{CO2,H2O}/\ce{HCO3-} (exercice~12).
\end{definition}
\begin{experience}[Déterminer un pKa par pH-métrie]
On mélange un volume d'une solution d'acide éthanoïque et un volume égal d'une solution d'éthanoate de sodium de même
concentration : alors $[\ce{AH}]_\text{éq}\approx[\ce{A-}]_\text{éq}$. Après étalonnage du pH-mètre, on mesure
$\pH=4{,}8$ à 25~°C : c'est le $\pKa$ du couple. En répétant la mesure pour d'autres proportions, on vérifie la relation
$\pH=\pKa+\log\left([\ce{A-}]/[\ce{AH}]\right)$.
\end{experience}
\begin{rediger}
Préciser toujours la température lorsqu'on utilise $K_\mathrm{e}$ ou un $\pKa$. Conclure une comparaison par une phrase :
« le $\pKa$ de l'acide méthanoïque est plus petit, cet acide est plus fort : à même concentration, son taux d'avancement
final est plus grand et son pH plus faible ». Vérifier qu'un résultat de pH est cohérent avec le diagramme de
prédominance.
\end{rediger}
\exercices

\rubrique{Découverte}
\begin{blocLivre}{ex-9-1}
\begin{exercice}[1]{Le produit ionique de l'eau}
À 25~°C, $K_\mathrm{e}=\num{1.0e-14}$ ; à 60~°C, $K_\mathrm{e}=\num{9.6e-14}$.
\begin{questions}
\item Écrire l'équation de la réaction d'autoprotolyse de l'eau et l'expression de $K_\mathrm{e}$.
\item Une boisson a un pH de 2,0 à 25~°C. Calculer $[\ce{H3O+}]$ puis $[\ce{HO-}]$.
\item Une eau savonneuse a un pH de 11,0 à 25~°C. Calculer $[\ce{HO-}]$.
\item Montrer que le pH d'une eau neutre vaut $\frac12\,\mathrm{p}K_\mathrm{e}$. Calculer le pH d'une eau pure à 60~°C.
Cette eau est-elle acide ?
\end{questions}
\end{exercice}
\end{blocLivre}
\begin{blocLivre}{ex-9-2}
\begin{exercice}[1]{Acide fort ou acide faible ?}
On dispose de deux solutions de même concentration apportée $c=\qty{1,0e-2}{mol.L^{-1}}$ : une solution d'acide
chlorhydrique, de pH 2,0, et une solution d'acide éthanoïque, de pH 3,4.
\begin{questions}
\item Définir un acide fort et un acide faible.
\item Écrire l'équation de la réaction de chaque acide avec l'eau, avec le symbole adapté.
\item Calculer le taux d'avancement final de chaque réaction. Conclure.
\item Laquelle des deux solutions contient le plus de molécules d'acide non dissociées ?
\end{questions}
\end{exercice}
\end{blocLivre}
\begin{blocLivre}{ex-9-3}
\begin{exercice}[1]{Écrire des constantes d'acidité}
On considère les couples \ce{HCOOH}/\ce{HCOO-}, \ce{NH4+}/\ce{NH3} et \ce{HClO}/\ce{ClO-}.
\begin{questions}
\item Pour chaque couple, écrire l'équation de la réaction de l'acide avec l'eau et l'expression de la constante
d'acidité.
\item Les $\pKa$ valent, dans le désordre : 9,2 ; 3,75 ; 7,5. Attribuer chaque valeur à son couple à l'aide du tableau
de l'essentiel, puis calculer les trois constantes $K_\mathrm{A}$.
\item Classer les trois acides du plus fort au plus faible.
\end{questions}
\end{exercice}
\end{blocLivre}
\begin{blocLivre}{ex-9-4}
\begin{exercice}[1]{Lire un diagramme de prédominance}
L'acide benzoïque \ce{C6H5COOH}, conservateur alimentaire, a un $\pKa$ de 4,2.
\schema{09-ex04}
\begin{questions}
\item Reproduire le diagramme de prédominance du couple et y placer les trois boissons.
\item Quelle forme du conservateur prédomine dans chaque boisson ?
\item Pour le soda, calculer le rapport $\dfrac{[\ce{C6H5COOH}]}{[\ce{C6H5COO-}]}$. Que peut-on dire de la forme
basique ?
\item Dans quelle boisson les deux formes sont-elles présentes en quantités égales ?
\end{questions}
\end{exercice}
\end{blocLivre}
\begin{blocLivre}{ex-9-5}
\begin{exercice}[1]{Du pH à la concentration}
Toutes les solutions sont à 25~°C.
\begin{questions}
\item Calculer le pH d'une solution d'acide nitrique (acide fort) de concentration \qty{2,0e-3}{mol.L^{-1}}.
\item Calculer le pH d'une solution de soude de concentration \qty{1,0e-3}{mol.L^{-1}}.
\item Un déboucheur de canalisations dilué a un pH de 10,4. Calculer la concentration en ions hydroxyde apportés.
\item Quel volume d'eau faut-il ajouter à 10,0~mL de la solution d'acide nitrique pour que son pH augmente d'une
unité ?
\end{questions}
\end{exercice}
\end{blocLivre}
\rubrique{Entraînement}
\begin{blocLivre}{ex-9-6}
\begin{exercice}[2]{Comparer la force de deux acides}
On prépare une solution d'acide méthanoïque et une solution d'acide éthanoïque, toutes deux à
$c=\qty{1,0e-2}{mol.L^{-1}}$.
\begin{donnees}
$\pKa(\ce{HCOOH}/\ce{HCOO-})=3{,}75$ ; $\pKa(\ce{CH3COOH}/\ce{CH3COO-})=4{,}76$.
\end{donnees}
\begin{questions}
\item Sans calcul, prévoir laquelle des deux solutions a le pH le plus faible. Justifier.
\item \capacite{Capacité mathématique}Pour l'acide méthanoïque, montrer que $h=[\ce{H3O+}]_\text{éq}$ vérifie
$h^2+K_\mathrm{A}h-K_\mathrm{A}c=0$, puis calculer $h$ et le pH.
\item Le pH de la solution d'acide éthanoïque vaut 3,39. Calculer le taux d'avancement final des deux réactions avec
l'eau. Le classement est-il cohérent avec les $\pKa$ ?
\item Les deux solutions sont diluées dix fois. Comment chaque taux d'avancement final évolue-t-il ?
\end{questions}
\end{exercice}
\end{blocLivre}
\begin{blocLivre}{ex-9-7}
\begin{exercice}[2]{Exploiter un diagramme de distribution}
Le graphique représente les proportions des deux formes du couple de l'acide éthanoïque en fonction du pH.
\schema{09-ex07}
\begin{questions}
\item Identifier la courbe associée à chaque forme. Justifier.
\item Déterminer graphiquement le $\pKa$ du couple.
\item Déterminer les proportions des deux formes à pH 3,8, puis le rapport $[\ce{CH3COO-}]/[\ce{CH3COOH}]$.
\item Retrouver ce rapport par le calcul à l'aide de la relation entre pH et $\pKa$.
\item Dans le vinaigre, de pH voisin de 2,8, sous quelle forme l'acide éthanoïque se trouve-t-il presque
exclusivement ?
\end{questions}
\end{exercice}
\end{blocLivre}
\begin{blocLivre}{ex-9-8}
\begin{exercice}[2]{Le pH d'une solution d'acide méthanoïque}
On dissout de l'acide méthanoïque dans de l'eau : la concentration apportée vaut $c=\qty{5,0e-3}{mol.L^{-1}}$.
\begin{donnees}
$K_\mathrm{A}=\num{1.8e-4}$ à 25~°C.
\end{donnees}
\begin{questions}
\item Dresser le tableau d'avancement de la réaction de l'acide avec l'eau, pour un volume $V$ de solution.
\item \capacite{Capacité mathématique}Établir l'équation du second degré vérifiée par $h=[\ce{H3O+}]_\text{éq}$, puis
la résoudre. En déduire le pH de la solution.
\item Calculer le taux d'avancement final.
\item On utilise parfois l'approximation $h\approx\sqrt{K_\mathrm{A}\,c}$, valable si l'acide est peu dissocié.
Calculer le pH approché et commenter l'écart.
\end{questions}
\end{exercice}
\end{blocLivre}
\begin{blocLivre}{ex-9-9}
\begin{exercice}[2]{Le bleu de bromothymol}
Le bleu de bromothymol (BBT) est un indicateur coloré : sa forme acide HIn est jaune, sa forme basique \ce{In-} est
bleue. Le $\pKa$ du couple vaut 7,1 ; la zone de virage s'étend de 6,0 à 7,6.
\schema{09-ex09}
\begin{questions}
\item Pourquoi un indicateur coloré change-t-il de couleur selon le pH ?
\item Prévoir la couleur du BBT dans le suc gastrique (pH 2), dans l'eau de mer (pH 8,2) et dans une eau pure à
25~°C.
\item Calculer le rapport $[\ce{In-}]/[\ce{HIn}]$ à pH 7,4.
\item Pourquoi la zone de virage est-elle large d'environ deux unités de pH ?
\item Quelques gouttes suffisent. Pourquoi faut-il éviter d'en verser davantage dans la solution étudiée ?
\end{questions}
\end{exercice}
\end{blocLivre}
\begin{blocLivre}{ex-9-10}
\begin{exercice}[2]{Une solution tampon}
On dissout 0,100~mol d'acide éthanoïque et 0,100~mol d'éthanoate de sodium dans de l'eau, de façon à obtenir 1,0~L de
solution S.
\begin{donnees}
$\pKa(\ce{CH3COOH}/\ce{CH3COO-})=4{,}76$ ; on admet que les quantités d'acide et de base conjuguée ne sont pas
modifiées de façon notable par leurs réactions avec l'eau.
\end{donnees}
\begin{questions}
\item Montrer que le pH de la solution S est égal au $\pKa$ du couple.
\item On ajoute \qty{1,0e-3}{mol} de chlorure d'hydrogène, qui transforme une partie des ions éthanoate en acide
éthanoïque. Déterminer les nouvelles quantités, puis le nouveau pH.
\item On ajoute la même quantité de chlorure d'hydrogène dans 1,0~L d'eau pure. Calculer le pH obtenu.
\item Comparer les deux variations de pH. Comment qualifie-t-on la solution S ?
\item Citer une situation où un tel contrôle du pH est vital.
\end{questions}
\end{exercice}
\end{blocLivre}
\begin{blocLivre}{ex-9-11}
\begin{exercice}[2]{Tracer un diagramme de distribution}
Le programme suivant calcule la proportion $\alpha$ de la forme basique d'un couple de $\pKa$ donné, pour trois valeurs
du pH.
\begin{python}
pKa = 4.76
def alpha(pH):
    return 1 / (1 + 10**(pKa - pH))
for pH in (3.76, 4.76, 5.76):
    print(round(alpha(pH), 3))
\end{python}
\begin{console}
0.091
0.5
0.909
\end{console}
\begin{questions}
\item \capacite{Capacité numérique}Établir l'expression de $\alpha$ utilisée par la fonction, à partir de la relation
entre pH et $\pKa$.
\item Interpréter les trois valeurs affichées.
\item Que vaut la proportion de la forme acide à $\pH=\pKa+1$ ? Quelle symétrie le diagramme de distribution
présente-t-il ?
\item Modifier le programme pour qu'il affiche la proportion de la forme acide à pH 7,4 pour le couple
\ce{CO2,H2O}/\ce{HCO3-}, de $\pKa$ 6,1 à 37~°C. Quelle valeur obtient-on ?
\end{questions}
\end{exercice}
\end{blocLivre}
\rubrique{Objectif bac}
\begin{blocLivre}{ex-9-12}
\begin{exercice}[3]{Le pH du sang}
\begin{docu}{1}{Une constante vitale}
Le pH du sang artériel est maintenu entre 7,35 et 7,45. Le principal système tampon du sang est le couple
\ce{CO2,H2O}/\ce{HCO3-} : le dioxyde de carbone dissous, produit par la respiration cellulaire, est éliminé dans l'air
expiré, tandis que les reins régulent la concentration en ions hydrogénocarbonate.
\end{docu}
\begin{docu}{2}{Effort et hyperventilation}
Lors d'un effort intense, la production de dioxyde de carbone et d'acide lactique augmente : le pH du sang diminue
légèrement, ce qui accélère la ventilation. À l'inverse, une hyperventilation prolongée élimine trop de dioxyde de
carbone et fait monter le pH : c'est l'alcalose respiratoire.
\end{docu}
\begin{donnees}
$\pKa(\ce{CO2,H2O}/\ce{HCO3-})=6{,}1$ à 37~°C ; dans le sang, $[\ce{HCO3-}]\approx\qty{24}{mmol.L^{-1}}$.
\end{donnees}
\begin{questions}
\item Écrire l'équation de la réaction du dioxyde de carbone dissous, noté \ce{CO2,H2O}, avec l'eau, qui forme des ions
hydrogénocarbonate.
\item Quelle forme du couple prédomine à pH 7,4 ? Justifier à l'aide d'un diagramme de prédominance.
\item \capacite{Capacité mathématique}Calculer le rapport $[\ce{HCO3-}]/[\ce{CO2,H2O}]$ à pH 7,4, puis la concentration
en dioxyde de carbone dissous.
\item Expliquer pourquoi un effort intense fait baisser le pH du sang, et pourquoi l'organisme réagit en accélérant la
ventilation.
\item \initiative Montrer qu'une diminution de moitié de $[\ce{CO2,H2O}]$, à $[\ce{HCO3-}]$ constante, fait monter le
pH d'environ 0,3 unité. Quel document ce calcul illustre-t-il ?
\end{questions}
\end{exercice}
\end{blocLivre}
\begin{blocLivre}{ex-9-13}
\begin{exercice}[3]{Le degré Dornic d'un lait}
\begin{docu}{1}{L'acidité d'un lait}
La fermentation du lactose produit de l'acide lactique \ce{C3H6O3} : plus un lait vieillit, plus il est acide.
L'acidité s'exprime en degrés Dornic : 1~°D correspond à 0,10~g d'acide lactique par litre. Un lait frais titre moins
de 18~°D ; au-delà de 25~°D, il caille à l'ébullition.
\end{docu}
\begin{donnees}
$M(\ce{C3H6O3})=\qty{90,1}{g.mol^{-1}}$ ; $\pKa(\text{acide lactique/ion lactate})=3{,}86$, soit
$K_\mathrm{A}=\num{1.4e-4}$.
\end{donnees}
\begin{questions}
\item Un lait titre 16~°D. Ce lait est-il frais ? Calculer la concentration molaire en acide lactique correspondante.
\item \capacite{Capacité mathématique}Calculer le pH d'une solution d'acide lactique de même concentration, puis son
taux d'avancement final.
\item Le pH du lait vaut en réalité 6,7. Sous quelle forme du couple l'acide lactique s'y trouve-t-il
majoritairement ? Proposer une explication à l'écart entre les deux pH.
\item \initiative Dans un yaourt, le pH vaut 4,5. Calculer le rapport entre les deux formes du couple. Le résultat
est-il cohérent avec le goût acide du yaourt ?
\end{questions}
\end{exercice}
\end{blocLivre}
\begin{blocLivre}{ex-9-14}
\begin{exercice}[3]{Un conservateur et son efficacité}
\begin{docu}{1}{L'acide benzoïque, conservateur E210}
L'acide benzoïque \ce{C6H5COOH} et les benzoates (E210 à E213) empêchent le développement des levures et des
moisissures dans les boissons. Seule la forme acide \ce{C6H5COOH}, capable de traverser la membrane des
micro-organismes, est réellement active. La réglementation européenne limite sa concentration à 150~mg par litre dans
les boissons non alcoolisées.
\end{docu}
\begin{donnees}
$\pKa(\ce{C6H5COOH}/\ce{C6H5COO-})=4{,}20$ ; $M(\ce{C6H5COOH})=\qty{122,1}{g.mol^{-1}}$ ; solubilité de l'acide
benzoïque dans l'eau à 25~°C : 2,9~g$\cdot$L$^{-1}$.
\end{donnees}
\begin{questions}
\item Calculer la concentration molaire maximale autorisée en conservateur. La solubilité permet-elle de l'atteindre ?
\item Dans un soda de pH 3,2, calculer la fraction du conservateur sous forme active.
\item Même question dans un thé glacé de pH 5,5. Comparer l'efficacité du conservateur dans les deux boissons.
\item \initiative Un fabricant souhaite conserver une boisson de pH 6,0 avec ce conservateur. Que lui
conseilleriez-vous ? Justifier par un calcul.
\end{questions}
\end{exercice}
\end{blocLivre}
\begin{blocLivre}{ex-9-15}
\begin{exercice}[3]{Désinfecter l'eau d'une piscine}
\begin{docu}{1}{Le chlore des piscines}
Le « chlore » dissous dans l'eau d'une piscine se répartit entre l'acide hypochloreux \ce{HClO} et l'ion hypochlorite
\ce{ClO-}. Seul l'acide hypochloreux, qui pénètre dans les micro-organismes, est fortement désinfectant : il est
environ cent fois plus efficace que l'ion hypochlorite.
\end{docu}
\begin{docu}{2}{Le réglage du pH}
Les piscinistes recommandent de maintenir le pH entre 7,2 et 7,4 : ce compromis préserve l'efficacité du chlore tout
en restant confortable pour la peau et les yeux, dont le pH de surface est proche de 7.
\end{docu}
\schema{09-ex15}
\begin{donnees}
$\pKa(\ce{HClO}/\ce{ClO-})=7{,}5$ à 25~°C.
\end{donnees}
\begin{questions}
\item Écrire la réaction de l'acide hypochloreux avec l'eau et l'expression de sa constante d'acidité.
\item À l'aide du graphique, déterminer la fraction d'acide hypochloreux à pH 7,2, puis à pH 8,2.
\item Retrouver ces deux valeurs par le calcul.
\item Un baigneur constate que l'eau de sa piscine, à pH 8,2, verdit malgré un apport correct de chlore. Expliquer, puis
indiquer comment rétablir la désinfection sans ajouter de chlore.
\item \initiative À l'aide des deux documents, justifier la plage de pH recommandée.
\end{questions}
\end{exercice}
\end{blocLivre}
